%============================================================
% CH1 · PATTERN 3 — REVERSAL AND ROTATION: THEORY
%============================================================
\subsection{Pattern Theory}
\label{ssec:ch1-p3-theory}

\begin{patternbox}{Pattern 3 at a Glance}
Two indices start at the two ends of a range and walk towards each other, swapping as they go: that \textbf{reverses} the range in place. Three reversals then \textbf{rotate} the whole array, again in place and in linear time.

\begin{center}
\begin{tikzpicture}[x=1cm, y=1cm]
    \dparray{a}{2.4}{start}{1/acell, 2/acell, 3/acell, 4/acell, 5/aread, 6/aread, 7/aread}
    \dparray{b}{1.4}{reverse all}{7/awrite, 6/awrite, 5/awrite, 4/awrite, 3/awrite, 2/awrite, 1/awrite}
    \dparray{c}{0.4}{reverse first $k$}{5/akeep, 6/akeep, 7/akeep, 4/acell, 3/acell, 2/acell, 1/acell}
    \dparray{d}{-0.6}{reverse the rest}{5/akeep, 6/akeep, 7/akeep, 1/akeep, 2/akeep, 3/akeep, 4/akeep}
    \node[note, anchor=west, text width=4.2cm, align=left] at (6.8,0.9)
        {$n = 7$, rotate right by $k = 3$: the last three values (orange) move to the front, the order inside each block is kept};
\end{tikzpicture}
\end{center}

\begin{recurrencebox}
\centering
\texttt{reverse(l, r)}:\quad \textbf{while} $l < r-1$:\; swap$(a[l], a[r-1])$,\; $l{+}{+}$,\; $r{-}{-}$ \qquad (range $[l, r)$)\\[3pt]
\texttt{rotateRight(k)}:\quad $k \leftarrow k \bmod n$;\; reverse$(0, n)$;\; reverse$(0, k)$;\; reverse$(k, n)$\\[2pt]
$\bigO{n}$ time, $\bigO{1}$ extra memory; every element is moved at most twice
\end{recurrencebox}
The same toolbox answers ``rotate left'', ``rotate a part of the array'', ``is $b$ a rotation of $a$?'' and ``was this array sorted and then rotated?''.
\end{patternbox}

%------------------------------------------------------------
\subsubsection*{From Pattern 2 to Pattern 3}

In Pattern~2 (Section~\ref{sec:ch1-p2}) two indices moved in the \emph{same} direction: a read index scanned every element and a write index trailed behind it, so the array was compacted without a second buffer. Pattern~3 keeps the goal --- rewrite the array in place, $\bigO{1}$ extra memory --- but sends the two indices in \emph{opposite} directions. The left index starts at the first cell, the right one at the last cell, and each step swaps the two cells they point at. Lesson~\ref{ssec:ch1-inplace} listed reversal as one of the basic in-place moves; here it becomes the building block of rotation.

\begin{diagrambox}
\begin{tikzpicture}[x=1cm, y=1cm]
    \begin{scope}
        \node[font=\small\bfseries, text=deepblue, anchor=west] at (-0.6,1.95) {Pattern 2: same direction};
        \dparray{p}{0.3}{}{1/akeep, 1/akeep, 2/akeep, 2/aout, 3/aread, 3/acell}
        \aptr{p3}{w}{accentorange}
        \aptr{p4}{r}{successgreen}
        \draw[flow] ([yshift=4mm]p4.north) to[bend right=40] node[lbl, above]{copy} ([yshift=4mm]p3.north);
    \end{scope}
    \begin{scope}[xshift=7.2cm]
        \node[font=\small\bfseries, text=deepblue, anchor=west] at (-0.6,1.95) {Pattern 3: towards each other};
        \dparray{q}{0.3}{}{6/akeep, 2/aread, 3/acell, 4/acell, 5/aread, 1/akeep}
        \aptr{q1}{l}{successgreen}
        \aptr{q4}{r}{accentorange}
        \draw[dep, <->] ([yshift=4mm]q1.north) to[bend left=35] node[lbl, above]{swap} ([yshift=4mm]q4.north);
    \end{scope}
\end{tikzpicture}
\figcap{Left: a read and a write index chase each other (compaction). Right: two indices meet in the middle (reversal); the green cells are already final.}
\end{diagrambox}

%------------------------------------------------------------
\subsubsection*{When to recognise it --- and when not}

Use Pattern 3 when
\begin{itemize}[leftmargin=*]
    \item the statement says \textbf{reverse}, \textbf{rotate}, \textbf{cyclic shift}, ``move the last $k$ elements to the front'', ``shift every element $k$ places to the right, wrapping around'';
    \item the order inside two blocks must be kept while the blocks \textbf{swap places} ($AB \to BA$): that is a rotation, whatever the story says;
    \item you must decide whether an array is a \textbf{rotation} of another one, or of a sorted array (Problem~\ref{ssec:ch1-sorted-rotated});
    \item many rotations are applied and only the final array (or a few cells) is needed: keep a \textbf{start offset} instead of moving data;
    \item the follow-up asks for $\bigO{1}$ extra memory.
\end{itemize}

\textbf{Do NOT use it} when the elements must be \emph{reordered by value} (that is sorting, Chapter~\ref{ch:sorting}), when a rotation is applied to a \emph{changing} subarray many times (implicit treaps, \treebook{27}{Treaps \& Cartesian Trees}), or when the question is a \emph{search} in a rotated sorted array --- that is binary search (Chapter~\ref{ch:binary_search}). Rotations of whole circular arrays come back in Chapter~\ref{ch:circular}.

%------------------------------------------------------------
\subsubsection*{Tool 1 --- reversing a range with two indices}

To reverse the half-open range $[l, r)$ keep $i = l$ and $j = r-1$; while $i < j$ swap $a[i]$ and $a[j]$, then move both inwards.

\begin{diagrambox}
\begin{tikzpicture}[x=1cm, y=1cm]
    \dparray{s}{2.0}{step 0}{a/aread, b/acell, c/acell, d/acell, e/acell, f/aread}
    \dparray{t}{0.6}{step 1}{f/akeep, b/aread, c/acell, d/acell, e/aread, a/akeep}
    \dparray{u}{-0.8}{step 2}{f/akeep, e/akeep, c/aread, d/aread, b/akeep, a/akeep}
    \dparray{v}{-2.2}{done}{f/adone, e/adone, d/adone, c/adone, b/adone, a/adone}
    \foreach \r/\x/\y in {s/0/5, t/1/4, u/2/3}{
        \draw[dep, <->] (\r\x.north) to[bend left=28] (\r\y.north);
    }
    \node[note, anchor=west, text width=5cm, align=left] at (5.9,0)
        {$n = 6$: three swaps, then $i = 3 > j = 2$ and the loop stops.\\[3pt]
        With $n = 7$ the middle cell is never touched --- it is its own mirror.\\[3pt]
        In general: $\lfloor n/2 \rfloor$ swaps.};
\end{tikzpicture}
\figcap{Reversing $[0,6)$ in place. After step $t$ the outer $t$ cells on each side are final (green).}
\end{diagrambox}

\begin{invariantbox}[Invariant of the reversal loop]
Before every test of $i < j$: the cells outside $[i, j]$ already hold their final (mirrored) values, $a[l+t] = a_{\text{old}}[r-1-t]$ for every $t < i-l$ and symmetrically on the right; the cells inside $[i, j]$ are untouched.

\textbf{Proof.} Initially nothing is outside. A swap puts $a_{\text{old}}[j]$ into position $i$ and $a_{\text{old}}[i]$ into position $j$ --- exactly the mirror positions, since $i - l = (r-1) - j$ is kept by moving both indices by one. When $i \ge j$ at most one cell (the middle) is left inside, and it is its own mirror. \hfill$\blacksquare$
\end{invariantbox}

%------------------------------------------------------------
\subsubsection*{Tool 2 --- rotation by three reversals}

\textbf{Rotating right by $k$} means that the element at index $i$ moves to index $(i + k) \bmod n$; the last $k$ elements come to the front. Write the array as two blocks: $A$ = the first $n-k$ elements, $B$ = the last $k$. We want $AB \to BA$. Let $X^R$ denote $X$ reversed. Reversal turns a concatenation around \emph{and} reverses each part:
\[
(AB)^R = B^R A^R, \qquad (X^R)^R = X .
\]
So reverse the whole array to get $B^R A^R$; now the blocks are in the right places but each one is backwards. Reverse the first $k$ cells ($B^R \to B$) and the last $n-k$ cells ($A^R \to A$):
\[
AB \;\xrightarrow{\ \text{reverse all}\ }\; B^R A^R \;\xrightarrow{\ \text{reverse }[0,k)\ }\; B\,A^R \;\xrightarrow{\ \text{reverse }[k,n)\ }\; BA .
\]

\begin{diagrambox}
\begin{tikzpicture}[x=1cm, y=1cm,
    blkA/.style={rectangle, draw=primaryblue, fill=primaryblue!15, minimum height=7mm, font=\small\bfseries},
    blkB/.style={rectangle, draw=accentorange, fill=answerfill, minimum height=7mm, font=\small\bfseries}]
    \node[blkA, minimum width=5cm, anchor=west] at (0,3) {$A$ \;$\longrightarrow$};
    \node[blkB, minimum width=3cm, anchor=west] at (5,3) {$B$ \;$\longrightarrow$};
    \node[blkB, minimum width=3cm, anchor=west] at (0,1.9) {$\longleftarrow$\; $B^R$};
    \node[blkA, minimum width=5cm, anchor=west] at (3,1.9) {$\longleftarrow$\; $A^R$};
    \node[blkB, minimum width=3cm, anchor=west] at (0,0.8) {$B$ \;$\longrightarrow$};
    \node[blkA, minimum width=5cm, anchor=west] at (3,0.8) {$\longleftarrow$\; $A^R$};
    \node[blkB, minimum width=3cm, anchor=west] at (0,-0.3) {$B$ \;$\longrightarrow$};
    \node[blkA, minimum width=5cm, anchor=west] at (3,-0.3) {$A$ \;$\longrightarrow$};
    \foreach \y/\t in {3/{start: $n-k$ then $k$}, 1.9/{reverse all}, 0.8/{reverse $[0,k)$}, -0.3/{reverse $[k,n)$}}{
        \node[font=\scriptsize\bfseries, text=deepblue, anchor=west] at (8.3,\y) {\t};
    }
\end{tikzpicture}
\figcap{The proof in pictures: arrows show the reading direction inside each block. Reversing everything swaps the blocks but flips them; two more reversals flip them back.}
\end{diagrambox}

\textbf{Rotating left by $k$} is rotating right by $n-k$; with reversals it is the same three calls in the other order: reverse $[0,k)$, reverse $[k,n)$, then reverse all ($AB$ with $|A| = k$ becomes $A^R B^R$, then $(A^R B^R)^R = BA$). Always reduce $k \bmod n$ first: rotating by $n$ is the identity, and $k$ may be much larger than $n$.

\textbf{Cost.} The three reversals do $\lfloor n/2\rfloor + \lfloor k/2\rfloor + \lfloor (n-k)/2\rfloor \le n$ swaps: $\bigO{n}$ time, $\bigO{1}$ extra memory. The standard library has \texttt{std::reverse(first, last)} and \texttt{std::rotate(first, middle, last)} (the latter makes \texttt{middle} the new first element, i.e.\ a \emph{left} rotation by \texttt{middle - first}).

%------------------------------------------------------------
\subsubsection*{Tool 3 --- the cycle (juggling) method and its $\gcd(n,k)$ cycles}

Rotation is a permutation of the positions: the value at $i$ must go to $p(i) = (i+k) \bmod n$. Follow it: pick up the value at position $s$, drop it at $p(s)$, pick up what was there, drop it at $p(p(s))$, and so on until you return to $s$. Every value is moved \emph{exactly once}, with one spare variable.

How many cycles are there? Starting at $s$ we visit $s, s+k, s+2k, \dots \pmod n$; we are back at $s$ after $t$ steps exactly when $n \mid tk$, i.e.\ after $t = n / \gcd(n,k)$ steps. So every cycle has length $n/g$ with $g = \gcd(n,k)$, there are exactly $g$ cycles, and the cycle of $s$ contains exactly the positions $\equiv s \pmod g$. Hence the starts $0, 1, \dots, g-1$ cover every position once.

\begin{diagrambox}
\begin{tikzpicture}[x=1cm, y=1cm]
    \begin{scope}
        \foreach \i in {0,...,5}{
            \pgfmathsetmacro{\ang}{90-60*\i}
            \ifodd\i \node[gvB] (c\i) at (\ang:1.35) {\i}; \else \node[gvA] (c\i) at (\ang:1.35) {\i}; \fi
        }
        \draw[gde, draw=primaryblue, line width=1.2pt] (c0) -- (c2); \draw[gde, draw=primaryblue, line width=1.2pt] (c2) -- (c4);
        \draw[gde, draw=primaryblue, line width=1.2pt] (c4) -- (c0);
        \draw[gde, draw=recursionred, line width=1.2pt] (c1) -- (c3); \draw[gde, draw=recursionred, line width=1.2pt] (c3) -- (c5);
        \draw[gde, draw=recursionred, line width=1.2pt] (c5) -- (c1);
        \node[font=\small\bfseries] at (0,-2.25) {$n = 6,\ k = 2$: $\gcd = 2$ cycles of length 3};
    \end{scope}
    \begin{scope}[xshift=6.6cm]
        \foreach \i in {0,...,6}{
            \pgfmathsetmacro{\ang}{90-360/7*\i}
            \node[gvA] (d\i) at (\ang:1.35) {\i};
        }
        \foreach \i/\j in {0/3, 3/6, 6/2, 2/5, 5/1, 1/4, 4/0}{
            \draw[gde, draw=primaryblue, line width=1.1pt] (d\i) -- (d\j);
        }
        \node[font=\small\bfseries] at (0,-2.25) {$n = 7,\ k = 3$: one cycle through all 7};
    \end{scope}
\end{tikzpicture}
\figcap{Arrows $i \to (i+k) \bmod n$. Left: positions split by parity into two cycles, so two starting points (0 and 1) are needed. Right: $\gcd(7,3) = 1$, a single cycle.}
\end{diagrambox}

On $[1,2,3,4,5,6]$ with $k = 2$ (running the template below): the cycle $0 \to 2 \to 4 \to 0$ gives $[5,2,1,4,3,6]$, then the cycle $1 \to 3 \to 5 \to 1$ gives $[5,6,1,2,3,4]$. With $k = 4$ the cycles are $0 \to 4 \to 2 \to 0$ and $1 \to 5 \to 3 \to 1$ and the result is $[3,4,5,6,1,2]$.

\begin{insightbox}[Reversal or cycles?]
Both are $\bigO{n}$ time and $\bigO{1}$ memory. Reversal does about $n$ swaps ($2n$ writes) with perfectly sequential memory access; the cycle method does exactly $n$ writes but jumps by $k$ cells, which is unfriendly to the cache for large arrays. Reversal is also much harder to get wrong. Use reversal by default; the cycle idea matters because the same ``follow the permutation'' loop applies \emph{any} permutation in place (Problem~\ref{ssec:ch1-build-permutation}).
\end{insightbox}

%------------------------------------------------------------
\subsubsection*{Tool 4 --- rotation without moving anything: a start offset}

Often the array does not need to be rotated at all. After rotating right by $k$, the element now at index $i$ is the one that was at index $(i - k) \bmod n$. Keep a single integer \texttt{start} and read through it:
\[
\text{rotated}[i] \;=\; a[(\texttt{start} + i) \bmod n], \qquad \text{rotate right by } k:\ \texttt{start} \leftarrow (\texttt{start} - k) \bmod n .
\]
A rotation then costs $\bigO{1}$, and $10^5$ rotations followed by a few lookups cost almost nothing. In C++ write \texttt{((start - k) \% n + n) \% n}: the \texttt{\%} of a negative number is negative.

\begin{diagrambox}
\begin{tikzpicture}[x=1cm, y=1cm]
    \dparray{m}{1.2}{stored $a$}{1/acell, 2/acell, 3/acell, 4/aread, 5/acell, 6/acell, 7/acell}
    \dpindex{mi}{1.95}{index}{0}{6}
    \aptr{m4}{start = 4}{accentorange}
    \dparray{v}{-0.9}{seen as}{5/akeep, 6/akeep, 7/akeep, 1/akeep, 2/akeep, 3/akeep, 4/akeep}
    \node[note, anchor=west, text width=4.3cm, align=left] at (6.7,0.2)
        {after ``rotate right by 3'': $\texttt{start} = (0 - 3) \bmod 7 = 4$; the view starts at $a[4] = 5$ and wraps around.\\ Nothing was moved.};
\end{tikzpicture}
\figcap{The rotated array is a different \emph{reading} of the same memory: index $i$ of the view is cell $(4 + i) \bmod 7$.}
\end{diagrambox}

%------------------------------------------------------------
\subsubsection*{Tool 5 --- detecting rotations}

\textbf{Is $b$ a rotation of $a$?} Every rotation of $a$ is a window of length $n$ of the doubled array $a\,a$: the rotation that starts at index $s$ is $a[s..n-1]\,a[0..s-1] = (aa)[s..s+n-1]$. So $b$ is a rotation of $a$ iff $|a| = |b|$ and $b$ occurs in $aa$. Trying every start costs $\bigO{n^2}$ --- fine for $n \le 10^3$; linear-time matching (\stringbook{8}{The KMP Algorithm and the Prefix Function}) makes it $\bigO{n}$. In code we never build $aa$: index it as $a[j \bmod n]$.

\textbf{Is $a$ a sorted array rotated?} A non-decreasing array has no \emph{drop} $a[i] > a[i+1]$. Rotating it creates at most one drop --- at the seam --- if we also look at the circular pair $(a[n-1], a[0])$. Conversely, at most one circular drop means that reading from the element after the drop gives a non-decreasing sequence. That is Problem~\ref{ssec:ch1-sorted-rotated}.

\begin{diagrambox}
\begin{tikzpicture}[x=1cm, y=1cm]
    \dparray{w}{0.5}{$a\,a$}{3/acell, 4/acell, 5/acell, 1/awin, 2/awin, 3/awin, 4/awin, 5/awin, 1/acell, 2/acell}
    \arange{w3}{w7}{$b = [1,2,3,4,5]$ found at $s = 3$}{memoorange}
    \draw[gcut] ([xshift=0.45cm, yshift=4mm]w4.north) -- ([xshift=0.45cm, yshift=-4mm]w4.south);
    \node[aptrlbl, text=warningred, above] at ([xshift=0.45cm, yshift=4mm]w4.north) {copy 2 starts};
    \node[note, anchor=west, text width=3.3cm, align=left] at (9.2,0.5) {$a = [3,4,5,1,2]$: the sorted window starts right after the only drop $5 > 1$};
\end{tikzpicture}
\figcap{Rotations of $a$ are exactly the length-$n$ windows of $aa$; $a$ is a rotated sorted array iff one of these windows is sorted.}
\end{diagrambox}

%------------------------------------------------------------
\subsubsection*{Variations map}

\begin{diagrambox}
\begin{tikzpicture}[
    var/.style={rectangle, rounded corners=2pt, draw=patternborder, fill=patternbg,
                font=\scriptsize, align=center, text width=3.25cm, minimum height=1.25cm, inner sep=3pt}]
    \node[var, fill=answerfill, draw=accentorange, text width=5.2cm, font=\scriptsize\bfseries] (root) at (0,0)
        {Reversal \& rotation\\two indices meet; $AB \to BA$ by three reversals};
    \node[var] (v1) at (-5.4,-2) {\textbf{rotate the array}\\right / left by $k$, $k \gg n$ (Problem~\ref{ssec:ch1-rotate-array})};
    \node[var] (v2) at (-1.8,-2) {\textbf{rotated sorted?}\\count circular drops (Problem~\ref{ssec:ch1-sorted-rotated})};
    \node[var] (v3) at (1.8,-2)  {\textbf{reverse pieces}\\blocks of $k$, one segment (E3.1, E3.6)};
    \node[var] (v4) at (5.4,-2)  {\textbf{rotation as offset}\\$(i + k) \bmod n$, many shifts (E3.2)};
    \node[var] (v5) at (-3.6,-3.6) {\textbf{best rotation}\\update a score from one rotation to the next (E3.4)};
    \node[var] (v6) at (0,-3.6)    {\textbf{rotate a subarray}\\choose the segment and the shift (E3.3)};
    \node[var] (v7) at (3.6,-3.6)  {\textbf{what reversals can reach}\\any permutation $\Rightarrow$ compare multisets (E3.5)};
    \node[draw=patternborder, dashed, rounded corners=4pt, fit=(v1)(v4)(v5)(v7), inner sep=5pt] (grp) {};
    \draw[flow] (root.south) -- (grp.north);
\end{tikzpicture}
\figcap{Every box uses reversal, a rotation offset, or the drop count. Rotations of 2D grids and matrices return in Chapter~\ref{ch:matrices}.}
\end{diagrambox}

\textbf{Twists and how to handle them}
\begin{itemize}[leftmargin=*]
    \item \textbf{$k \ge n$ or $k = 0$:} reduce $k \bmod n$ first; after that $k = 0$ must still work (all three reversals are then harmless).
    \item \textbf{Left versus right:} left by $k$ = right by $n - k$. Check one tiny example by hand before coding.
    \item \textbf{Negative offsets:} $(s - k) \bmod n$ in C++ needs \texttt{((s - k) \% n + n) \% n}.
    \item \textbf{Rotate only $[l, r)$:} the same three reversals inside the range, with $k \bmod (r - l)$.
    \item \textbf{Duplicates:} ``rotated sorted'' tests must use $>$ (a strict drop), never $\ge$; equal neighbours are allowed in a non-decreasing array.
\end{itemize}

%------------------------------------------------------------
\subsubsection*{Code templates}

\begin{lstlisting}[style=cpp]
#include <bits/stdc++.h>
using namespace std;

// reverse the half-open range [l, r) in place
void reverseRange(vector<int>& a, int l, int r) {
    for (int i = l, j = r - 1; i < j; i++, j--) swap(a[i], a[j]);
}
// a[i] moves to (i + k) mod n
void rotateRight(vector<int>& a, long long k) {
    int n = a.size();
    if (n == 0) return;
    k %= n;
    reverseRange(a, 0, n);
    reverseRange(a, 0, k);
    reverseRange(a, k, n);
}
// a[i] moves to (i - k) mod n
void rotateLeft(vector<int>& a, long long k) {
    int n = a.size();
    if (n == 0) return;
    rotateRight(a, n - k % n);
}
// cycle (juggling) method: gcd(n, k) cycles, every value written once
void rotateRightCycles(vector<int>& a, int k) {
    int n = a.size();
    k %= n;
    if (k == 0) return;
    int g = __gcd(n, k);
    for (int s = 0; s < g; s++) {
        int cur = s, carry = a[s];
        do {
            cur = (cur + k) % n;
            swap(carry, a[cur]);
        } while (cur != s);
    }
}
// is b a rotation of a?  O(n^2) by trying every start in a a
bool isRotation(const vector<int>& a, const vector<int>& b) {
    int n = a.size();
    if ((int)b.size() != n) return false;
    if (n == 0) return true;
    for (int s = 0; s < n; s++) {
        int j = 0;
        while (j < n && a[(s + j) % n] == b[j]) j++;
        if (j == n) return true;
    }
    return false;
}

void print(const vector<int>& a) {
    for (size_t i = 0; i < a.size(); i++) cout << a[i] << (i + 1 < a.size() ? ' ' : '\n');
}

int main() {
    vector<int> a = {1, 2, 3, 4, 5, 6, 7};
    rotateRight(a, 3);          print(a);   // 5 6 7 1 2 3 4
    rotateLeft(a, 3);           print(a);   // 1 2 3 4 5 6 7
    rotateRight(a, 1000000007); print(a);   // k mod 7 = 6: 2 3 4 5 6 7 1
    vector<int> b = {1, 2, 3, 4, 5, 6};
    rotateRightCycles(b, 2);    print(b);   // 5 6 1 2 3 4
    cout << isRotation({3, 4, 5, 1, 2}, {1, 2, 3, 4, 5}) << " "
         << isRotation({1, 2, 3}, {1, 3, 2}) << "\n";   // 1 0
    // a rotation as an offset: view[i] = a[(start + i) % n]
    vector<int> c = {1, 2, 3, 4, 5, 6, 7};
    int n = c.size(), start = 0;
    for (int k : {3, 5, 2}) start = ((start - k) % n + n) % n;   // 3 rotations, total 10
    for (int i = 0; i < n; i++)                                  // 5 6 7 1 2 3 4
        cout << c[(start + i) % n] << (i + 1 < n ? ' ' : '\n');
}
\end{lstlisting}

\begin{insightbox}[Check the templates by hand]
Right by $10^9+7$ on $n = 7$ is right by $(10^9+7) \bmod 7 = 6$, i.e.\ left by 1: $[2,3,4,5,6,7,1]$. Three right rotations by $3, 5, 2$ are one rotation by $10 \equiv 3 \pmod 7$, hence the offset view prints $[5,6,7,1,2,3,4]$, the same as the first line. All six printed lines were produced by running the program.
\end{insightbox}

%------------------------------------------------------------
\subsubsection*{Common mistakes}

\begin{warningbox}[Mistake 1 --- forgetting $k \bmod n$]
\begin{lstlisting}[style=snippet]
reverse(a.begin(), a.begin() + k);      // WRONG when k > n: iterator past the end
k %= n; reverse(a.begin(), a.begin() + k);   // RIGHT
\end{lstlisting}
LeetCode~189 allows $k = 10^5$ with $n = 1$: without the reduction the program crashes or corrupts memory.
\end{warningbox}

\begin{warningbox}[Mistake 2 --- moving values in place in the obvious order]
\texttt{for i: a[(i+k)\%n] = a[i]} overwrites a value before it has been moved. On $[1,2,3,4,5,6,7]$ with $k=3$ it produces $[2,3,1,1,2,3,1]$. Either copy to a second array, or follow whole cycles with a spare variable (Tool~3), or use reversals.
\end{warningbox}

\begin{warningbox}[Mistake 3 --- the wrong order of the three reversals]
``reverse $[0,k)$, reverse $[k,n)$, reverse all'' rotates \emph{left} by $k$: $[1..7]$ becomes $[4,5,6,7,1,2,3]$ instead of $[5,6,7,1,2,3,4]$. Test with $n = 3$, $k = 1$ before submitting.
\end{warningbox}

\begin{warningbox}[Mistake 4 --- off-by-one in the reversal bounds]
With a half-open range the right index starts at $r-1$, not at $r$; the loop condition is $i < j$ (with $i \le j$ the middle cell is swapped with itself --- harmless, but a sign that the bounds are not under control).
\end{warningbox}

\begin{warningbox}[Mistake 5 --- linear drops instead of circular drops]
For ``sorted and rotated'' counting only the pairs $(i, i+1)$ with $i+1 < n$ misses the seam $(a[n-1], a[0])$: $[2,1,3,4]$ has one linear drop but is \emph{not} a rotated sorted array.
\end{warningbox}

\begin{checklistbox}[``I'm stuck'' checklist for Pattern 3]
\begin{itemize}[leftmargin=*]
    \item[$\square$] Can I describe the result as two blocks that swap places ($AB \to BA$)? Then it is a rotation.
    \item[$\square$] Right or left? What is $k \bmod n$? Did I test $k = 0$, $k = n$, $n = 1$?
    \item[$\square$] Do I need the rotated array, or only some cells of it? A start offset may be enough.
    \item[$\square$] Must I compare with \emph{all} rotations? Think of $aa$ and its windows of length $n$.
    \item[$\square$] Is the question about sortedness? Count the circular drops.
    \item[$\square$] Is the required memory $\bigO{1}$? Reversals, or cycles of the permutation.
\end{itemize}
\end{checklistbox}
